\section{Construction of Expression Vector Space}\label{sec:vecspace}
    \begin{itemize}
        \item Biological replicates belonging to the same axis (tissue, stressed tissue, etc.) are averaged post-normalization.
        \item Each gene is represented as a vector \( \vec{p} \in \mathbb{R}^n \),
              where \( n \) is the number of conditions or tissues.
              Hence the analysis is embedded in an $n$-dimensional ambient space.
        \item Expression vectors are stored in Fortran arrays. \emph{Importantly}
              Fortran is column major and vectors \emph{must} be columns.
        \item Two K-D Trees are constructed:
              \begin{itemize}
                  \item A raw-space K-D Tree (unnormalized vectors) for Euclidean queries.
                  \item A normalized-space K-D Tree (unit-length vectors) for cosine similarity and angular queries.
              \end{itemize}
        \item The first axis is designated the \textbf{Anchor Axis (AA)} for
              visualization.
    \end{itemize}

    \subsection{Vector Space Storage in Fortran}\label{sec:vecstore}

        All expression vectors in Tensor Omics are stored in a central memory structure optimized for performance and access in Fortran. This architecture ensures efficient storage, lookup, and downstream computation.

        \begin{itemize}
            \item \textbf{vec\_store:} A two-dimensional Fortran array of size \( n \times v \), where \( n \) is the number of expression dimensions (e.g., tissues or conditions), and \( v \) is the number of vectors (e.g., genes). Each column represents one expression vector.

            \item \textbf{cart\_tree:} A K-D Tree constructed on the raw (unnormalized) vectors stored in \texttt{vec\_store}. This structure enables rapid spatial neighborhood queries using Euclidean distance in the original expression space. This is not initialized until step~\ref{sec:kdtree}.

            \item \textbf{sphere\_tree:} A second K-D Tree constructed from unit-length normalized vectors derived from \texttt{vec\_store}. This structure supports directional queries based on cosine similarity or angular distance. This is not initialized until step~\ref{sec:kdtree}.
        \end{itemize}

        Together, these data structures allow both magnitude-sensitive and direction-sensitive operations across expression fields.